\documentclass[aps,pra,showpacs,notitlepage,onecolumn,superscriptaddress,nofootinbib]{revtex4-1}
\usepackage[utf8]{inputenc}
\usepackage[tmargin=1in, bmargin=1.25in, lmargin=1.5in, rmargin=1.5in]{geometry}
\usepackage{amsmath, amssymb, amsthm}
\usepackage{graphicx}
\usepackage{xcolor}
\usepackage{enumitem}
\usepackage{datetime}
\usepackage{hyperref}
\usepackage{titlesec}
\usepackage{import}
\usepackage{mathtools}
\usepackage{thmtools,thm-restate}
\usepackage{tikz-cd}
\usepackage{verbatim}


% package for commutative diagrams
% \usepackage{tikz-cd}

%%%%%%%%%%%%%%%%%%%%%%%%%%%%%%%%%%%%%%%%%%%%%
\definecolor{crimson}{RGB}{186,0,44}
\definecolor{moss}{RGB}{0, 186, 111}
\newcommand{\pop}[1]{\textcolor{crimson}{#1}}
\newcommand{\zcom}[1]{\noindent\textcolor{crimson}{(Z): #1}}
\newcommand{\jcom}[1]{\noindent\textcolor{moss}{(J): #1}}
\newcommand{\wt}[1]{\widetilde{#1}}
\newcommand{\pqeq}{\succcurlyeq}
\newcommand{\pleq}{\preccurlyeq}

%%%%%%%%%%%%%%%%%%%%%%%%%%%%%%%%%%%%%%%%%%%%%
\hypersetup{
    colorlinks,
    linkcolor={crimson},
    citecolor={crimson},
    urlcolor={crimson}
}

\usepackage{qcircuit}
\usepackage{comment}

%%%%%%%%%%%%%%%%%%%%%%%%%%%%%%%%%%%%%%%%%%%%%
\theoremstyle{definition}
\newtheorem{definition}{Definition}[section]
\newtheorem{lemma}{Lemma}[section]
\newtheorem{theorem}{Theorem}[section]
\newtheorem{corollary}{Corollary}[theorem]
\newtheorem*{theorem*}{Theorem}
\newtheorem*{corollary*}{Corollary}
\newtheorem{remark}{Remark}[section]
\newtheorem{conjecture}{Conjecture}[section]
\newtheorem{example}{Example}[section]
\newtheorem{reminder}{Reminder}[section]
\newtheorem{problem}{Problem}[section]
\newtheorem{question}{Question}[section]
\newtheorem{answer}{Answer}[section]
\newtheorem{fact}{Fact}[section]
\newtheorem{claim}{Claim}[section]
\newtheorem{proposition}{Proposition}[section]
\newtheorem{mantra}{Mantra}[section]

\usepackage{geometry}
\geometry{
  left=22mm,
  right=22mm,
  top=20mm,
}

\newcommand{\hhrulefill}{\hspace{-1.5em} \hrulefill}
\renewcommand{\baselinestretch}{1.1} 

%%%%%%%%%%%%%%%%%%%%%%%%%%%%%%%%%%%%%%%%%%%%%
\bibliographystyle{unsrt}

%%%%%%%%%%%%%%%%%%%%%%%%%%%%%%%%%%%%%%%%%%%%%
%%%%%%%%%%%%%%%%%%%%%%%%%%%%%%%%%%%%%%%%%%%%%
%%%%%%%%%%%%%%%%%%%%%%%%%%%%%%%%%%%%%%%%%%%%%
\begin{document}

\title{Hartshorne chapter 2 section 1 solutions}
\author{Jack Ceroni}
\date{\today}
\maketitle


\tableofcontents

\hhrulefill

\section{Introduction}

\noindent As a small disclaimer, note that throughout these solutions, equality is often taken to mean ``isomorphic via unique isomorphism''. In particular cases where it is convenient to make use of universal properties,
all instances of objects satisfying the universal property are canonically isomorphic (isomorphic via unique isomorphism), and thus for all intents and purposes, these objects may all be regarded
as the same.

\section{Solutions}

\subsection{Problem 2.1.1}

\noindent The constant presheaf $\mathcal{F}_p$ clearly has all of its stalks isomorphic to $A$. Therefore, some function $s : U \to \cup_{p \in U} \mathcal{F}_p$ with $s(p) \in \mathcal{F}_p$ and
for each $p \in U$ having $p \in V \subset U$ with $t \in \mathcal{F}(V) = A$ with $s(q) = t$ for all $q \in V$ is equivalent to specifying a function $s : U \to A$ which is continuous when $A$ is
given the discrete topology (in other words, a locally-constant function). The result immediately follows.

\subsection{Problem 2.1.2}

\noindent \textbf{Note:} I am going to use the notation $\text{im}(\varphi)$ for the image presheaf and $\text{im}(\varphi)^{+}$ for the corresponding sheafification.
\newline

\noindent \textbf{Part A.} The claim is that $\text{ker}(\varphi)_p$ and $\text{ker}(\varphi_p)$ are isomorphic as Abelian groups via a canonical isomorphism.
In particular, we have the subsheaf inclusion $j : \text{ker}(\varphi) \to \mathcal{F}$. Our goal is to show that the induced map on stalks $j_p : \text{ker}(\varphi)_p \to \mathcal{F}_p$
is an injective group homomorphism whose image is $\text{ker}(\varphi_p) \subset \mathcal{F}_p$. Since $j_p$ takes $s_p$ with $s \in \text{ker}(\varphi)(U) = \text{ker}(\varphi(U))$ to
$s_p$ with $s \in \mathcal{F}(U)$, injectivity is obvious and it is obvious that $j_p(\text{ker}(\varphi)_p) \subset \text{ker}(\varphi_p)$. On the other hand, if $\varphi_p(s_p) = 0$
for some $s_p \in \mathcal{F}_p$, then we must have $\varphi(W)(s|_W) = 0$ for some $W$ around $p$, so $s|_W \in \text{ker}(\varphi(W))$ and we have well-defined $(s|_W)_p \in \text{ker}(\varphi)_p$ where $j_p((s|_W)_p) = (s|_W)_p = s_p$.

In the case that $\mathcal{F}$ is a presheaf and $\mathcal{F}^{+}$ is the sheafification, recall that we have the canonical embedding $\theta : \mathcal{F} \to \mathcal{F}^{+}$
which is a presheaf morphism where $\theta(U)$ takes $s \in \mathcal{F}(U)$ to function $s^{+} : U \to \sqcup_{p \in U} \mathcal{F}_p$ such that $s^{+}(p) = s_p \in \mathcal{F}_p$.
It isn't difficult to see that $\theta_p : \mathcal{F}_p \to \mathcal{F}^{+}_p$ is an isomorphism on stalks where $\theta_p(s_p) = s_p^{+}$. Thus, if we wish to verify that the
subsheaf inclusion $j^{+} : \text{im}(\varphi)^{+} \to \mathcal{G}$ induces injective $j^{+}_p : \text{im}(\varphi)^{+}_p \to \mathcal{G}_p$ whose image is $\text{im}(\varphi_p) \subset \mathcal{G}_p$,
it suffices to show that this is true for the map on stalks induced by $j : \text{im}(\varphi) \to \mathcal{G}$, as $j = j^{+} \circ \theta$. The proof of this fact
is pretty much the same as the earlier proof for kernels.
\newline

\noindent \textbf{Part B.} $\varphi$ is injective if and only if $\text{ker}(\varphi) = 0$ which occurs if and only if $\text{ker}(\varphi)_p = 0$ for all $p$ which occurs if and only
if $\text{ker}(\varphi_p) = 0$ for all $p$ via Part A. In addition, given $\varphi : \mathcal{F} \to \mathcal{G}$, surjectivity is equivalent to showing that $j^{+} : \text{im}(\varphi)^{+} \to \mathcal{G}$
is an isomorphism, which is equivalent to showing that $j^{+}_p : \text{im}(\varphi)^{+}_p \to \mathcal{G}_p$ is an isomorphism of Abelian groups for each $p$. We already showed that this map
is injective with image $\text{im}(\varphi_p)$, hence it is an isomorphism if and only if $\text{im}(\varphi_p) = \mathcal{G}_p$ (i.e. $\varphi_p$ is surjective).
\newline

\noindent \textbf{Part C.} We have subsheaf inclusions $i : \text{ker}(\varphi^k) \to \mathcal{F}^k$ and $j^{+} : \text{im}(\varphi^{k-1})^{+} \to \mathcal{F}^k$. Exactness
is the requirement that there is unique isomorphism of sheaves $\phi : \text{ker}(\varphi^k) \to \text{im}(\varphi^{k-1})^{+}$ such that $i = j^{+} \circ \phi$. If such
an isomorphism exists, then on stalks $i_p = j^{+}_p \circ \phi_p$ and thus $\text{im}(i_p) = \text{im}(j^{+}_p)$, hence $\text{ker}(\varphi^k_p) = \text{im}(\varphi^{k-1}_p)$
from Part A, and we have exactness on stalks. Conversely, $\text{ker}(\varphi^k_p) = \text{im}(\varphi^{k-1}_p)$ implies that $\text{im}(i_p) = \text{im}(j^{+}_p)$. It follows that
for some $s \in \text{ker}(\varphi^k)(U)$ with $p \in U$, there is unique $t_p \in \text{im}(\varphi^{k-1})_p$ such that $i_p(s_p) = j^{+}_p(t^{+}_p) = j_p(t_p)$ in $\mathcal{F}^k_p$. In
particular, there must be open $W$ around $p$ such that $i(W)(s|_W) = j(W)(t|_W)$. This allows us to define $\phi(W)(s|_W) = \theta(W)(t|_W)$. The locality properties of sheaves allow
us to extend this to a well-defined morphism of sheaves, which is in fact an isomorphism, and is unique (as it induces an isomorphism on stalks, and is uniquely-determined on stalks).

\end{document}
